\section{Indledning}
\subsection{Baggrund}
I dag anvender det danske kajakpolofællesskab primært spreadsheets og manuelle processer til håndtering af
turneringsplaner, nedskrivning af resultater, pulje planlægning, match-making osv. og en primær
kommunikation via Facebook, hvilket resulterer i ineffektiv håndtering af dette. Det øger risikoen
for fejl, manglende overblik og begrænser tilgængeligheden for spillere, trænere og arrangører.

\subsection{Problemformulering}
Hvordan kan en webbaseret løsning udvikles til det danske kajakpolofællesskab, der understøtter
håndtering af turneringer, kampresultater, holdstatistikker m.m., med henblik på at reducere
afhængigheden af manuelle og spreadsheet-baserede processer og dermed skabe et mere effektivt, tilgængeligt
og ligetil system?

\subsection{Projektbeskrivelse}
Projektet har til formål at udvikle en webbaseret applikation, der effektivt understøtter
håndteringen af kajakpolo turneringer og aktiviteter i Danmark. Applikationen skal tilbyde funktioner til
Content Management på en offentlig front samt til administration af medlemskaber og kommunikation på klub-niveau. Funktionaliteterne vil inkludere
turneringsplanlægning, resultatregistrering, resultatstatistikker, holdstyring og kommunikationsværktøjer.
Ved at automatisere disse processer sigter projektet mod at reducere
fejl, forbedre overblikket og øge tilgængeligheden af information for alle involverede parter. Projektet vil
omfatte kravindsamling, design, implementering og testning af applikationen på en selvstyret server for at sikre, at den opfylder
brugernes behov og forventninger. Hovedfokus vil være på brugervenlighed, pålidelighed og skalerbarhed for at sikre en langvarig løsning for det danske kajakpolofællesskab.

\subsection{Afgrænsning}
Applikationens hovedmål er som beskrevet i projektbeskrivelsen. Til projektet vil der kun indgå en desktop-version
af applikationen, da det vurderes at være tilstrækkeligt for CMS- og administrationsfunktionaliteterne. Mobilversioner
og apps vil ikke blive udviklet i dette forløb. Ydermere vil der ikke blive implementeret betalingsløsninger eller integrationer med tredjepartstjenester.
Se bilag~\ref{bilag:afgrænsning} for en komplet liste over afgrænsninger.

\subsection{Interessenter}
Projektets primære interessenter omfatter:
\begin{itemize}
    \item \textbf{Dansk Kajakpolo Forbund (DKF)}
    \item \textbf{Individuelle kajakklubber- og forbund i Danmark}
    \item \textbf{Turneringsarrangører og Dommere}
    \item \textbf{Kajakpolo Spillere og Trænere tilknyttet klubberne}
\end{itemize}
Projektet er tilknyttet en representant fra Dansk Kajakpolo Forbund (MÅ vI SKRIVE HANS NAVN?), som fungerer som
projektets kunde og hovedinteressent. Denne person vil bidrage til kravindsamling, feedback og godkendelse af
projektets leverancer.